\chapter{O que o agente acerta e o que não acerta}
\label{cap:resultados-agente}

\begin{objetivos}
Ao final deste capítulo, você será capaz de:
\begin{itemize}[leftmargin=*,itemsep=0.2em]
  \item interpretar acurácia em três níveis de ordenação e explicar por que a diferença
        entre elas é informativa;
  \item avaliar um sistema no papel de validador, distinguindo precisão de revocação
        nesse contexto;
  \item analisar convergência de nomes gerados em linguagem natural;
  \item identificar artefatos de conjunto de dados a partir de padrões nas descrições
        produzidas;
  \item derivar uma divisão de trabalho entre modelos numéricos e camada de linguagem a
        partir de evidência empírica.
\end{itemize}
\end{objetivos}

\section{Três estudos, três perguntas}

O \cref{cap:agente} descreveu como o agente é construído. Este capítulo mede o que ele
faz, por meio de três estudos independentes que respondem a três perguntas distintas.

\begin{enumerate}[leftmargin=*,itemsep=0.3em]
  \item \textbf{Estudo 0 --- classificação.} O agente consegue nomear a classe correta a
        partir apenas das estatísticas?
  \item \textbf{Estudo 1 --- validação.} O agente consegue julgar se a decisão de um
        classificador numérico está correta?
  \item \textbf{Estudo 2 --- novidade.} O agente consegue reconhecer que um segmento não
        corresponde a nenhuma classe conhecida?
\end{enumerate}

Adianta-se a conclusão, porque ela organiza a leitura: \textbf{o agente é ruim na
primeira tarefa, razoável na segunda e bom na terceira}. E essa ordenação, longe de ser
um resultado decepcionante, é o que define seu papel correto no sistema.

\section{Estudo 0: classificação}

\subsection{Protocolo}

O agente recebe as métricas de um segmento e os perfis das nove classes, e produz três
candidatas ordenadas. Avaliam-se 191 segmentos de poços reais.

Reportam-se três níveis de acurácia: \emph{top-1} (a primeira candidata está correta),
\emph{top-2} (a correta está entre as duas primeiras) e \emph{top-3}. Intervalos de
confiança de $95\,\%$ são de Wilson (\cref{eq:wilson}).

\subsection{Resultados globais}

\begin{table}[htbp]
\centering
\caption{Acurácia global do agente na tarefa de classificação, com intervalos de
confiança de Wilson a $95\,\%$. As linhas de acaso correspondem a escolha aleatória
entre nove classes.}
\label{tab:ag-global}
\begin{tabular}{lccc}
\toprule
\textbf{Nível} & \textbf{Acurácia} & \textbf{IC $95\,\%$} & \textbf{Acaso} \\
\midrule
Top-1 & $35{,}1\,\%$ & $[28{,}7;\ 42{,}1]$ & $11{,}1\,\%$ \\
Top-2 & $53{,}4\,\%$ & $[46{,}3;\ 60{,}3]$ & $22{,}2\,\%$ \\
Top-3 & $63{,}9\,\%$ & $[56{,}9;\ 70{,}4]$ & $33{,}3\,\%$ \\
\bottomrule
\end{tabular}
\end{table}

\begin{resultadochave}
O agente acerta a classe na primeira tentativa em $35{,}1\,\%$ dos casos --- três vezes o
acaso --- e a inclui entre as três primeiras em $63{,}9\,\%$, quase o dobro do acaso.
Nenhum desses números é competitivo com os classificadores numéricos, que atingem
$0{,}87$ de acurácia (\cref{cap:binarios}).
\end{resultadochave}

\subsection{Como não ler estes números}

Duas leituras erradas são tentadoras, e ambas devem ser afastadas.

\textbf{``O agente é inútil.''} Ele opera em condições incomparavelmente mais duras:
recebe treze métricas por sensor em vez de duzentas amostras por sensor, nunca foi
treinado nesta tarefa, e não viu um único exemplo rotulado. Três vezes o acaso, sob essas
condições, indica que a linguagem intermediária do \cref{cap:agente} de fato carrega
informação diagnóstica.

\textbf{``O agente vai substituir os classificadores.''} $35{,}1\,\%$ contra $87\,\%$
encerra o assunto.

\begin{destaque}[A informação está na diferença entre os níveis]
De top-1 para top-3, a acurácia quase dobra: $35{,}1\,\%$ para $63{,}9\,\%$.

Isso significa que, em cerca de $29\,\%$ dos casos, o agente \emph{tem} a resposta certa
em mãos e a coloca em segundo ou terceiro lugar. Ele identifica o conjunto plausível de
hipóteses e falha em ordená-lo.

É exatamente o perfil de um sistema que deve ser usado como \textbf{gerador de
hipóteses}, e não como decisor. E é a justificativa empírica da decisão de projeto ---
tomada antes de conhecer estes resultados --- de exigir três candidatas em vez de uma.
\end{destaque}

\subsection{Por classe}

\begin{table}[htbp]
\centering
\caption{Acurácia do agente por classe na tarefa de classificação (\%).}
\label{tab:ag-porclasse}
\small
\begin{tabular}{llrccc}
\toprule
\textbf{Cl.} & \textbf{Evento} & $\bm{n}$ & \textbf{Top-1} & \textbf{Top-2} & \textbf{Top-3} \\
\midrule
1 & Aumento de BSW            & 15 & $26{,}7$ & $40{,}0$ & $53{,}3$ \\
2 & Fechamento de DHSV        & 22 & $\bm{86{,}4}$ & $\bm{90{,}9}$ & $\bm{100{,}0}$ \\
3 & Golfadas severas          & 32 & $18{,}8$ & $53{,}1$ & $53{,}1$ \\
4 & Instabilidade de fluxo    & 27 & $18{,}5$ & $25{,}9$ & $37{,}0$ \\
5 & Perda de produtividade    & 15 & $26{,}7$ & $53{,}3$ & $66{,}7$ \\
6 & Restrição rápida no PCK   & 15 & $66{,}7$ & $66{,}7$ & $66{,}7$ \\
7 & Incrustação no PCK        & 36 & $38{,}9$ & $66{,}7$ & $83{,}3$ \\
8 & Hidrato na produção       & 15 & $26{,}7$ & $53{,}3$ & $60{,}0$ \\
9 & Hidrato no serviço        & 14 & $\bm{7{,}1}$ & $\bm{14{,}3}$ & $42{,}9$ \\
\bottomrule
\end{tabular}
\end{table}

A dispersão é enorme --- de $7{,}1\,\%$ a $86{,}4\,\%$ na primeira escolha --- e os
extremos são instrutivos.

\paragraph{Classe 2, o melhor caso.} $86{,}4\,\%$ na primeira escolha e $100\,\%$ nas
três primeiras. O fechamento de DHSV tem assinatura inequívoca: queda abrupta e
simultânea de pressão a montante e a jusante, com direção e magnitude que nenhuma outra
classe reproduz. Quando a assinatura é distintiva, treze métricas bastam.

\paragraph{Classe 6, o caso rígido.} $66{,}7\,\%$ nos três níveis --- valor idêntico. O
agente ou acerta de primeira, ou não considera a classe correta em lugar algum. Não há
hesitação nem gradação. É o comportamento de um critério determinístico embutido no
perfil da classe, que ou é satisfeito ou não é.

\paragraph{Classe 9, o pior caso.} $7{,}1\,\%$ na primeira escolha, $42{,}9\,\%$ nas três
primeiras. Hidrato na linha de serviço é sistematicamente confundido com hidrato na linha
de produção (classe 8) --- e não poderia ser diferente: o mecanismo físico é o mesmo,
apenas a localização difere, e nenhuma das treze métricas codifica localização.

\begin{destaque}[Um erro que não é do modelo]
A confusão entre as classes 8 e 9 não é uma falha de raciocínio. É uma consequência
direta da escolha de métricas do \cref{cap:agente}, que foi apresentada lá como ``a
decisão de projeto mais consequente desta parte do livro''.

O que as métricas não capturam, o modelo não pode inferir. Distinguir 8 de 9 exigiria
informação de qual conjunto de sensores está afetado --- informação disponível, mas não
codificada.

\emph{Este é um problema corrigível.} Vários dos números baixos desta tabela podem estar
medindo a linguagem intermediária, não o agente.
\end{destaque}

\paragraph{Classes 4 e 5.} $18{,}5\,\%$ e $26{,}7\,\%$. Pela quarta vez neste livro, as
classes sintomáticas aparecem no fundo da tabela. A consistência do padrão --- em
métodos completamente distintos, de ensembles binários a modelos de linguagem --- é a
evidência mais forte de que o problema está na taxonomia.

Por essa razão, \textbf{as classes 4 e 5 foram excluídas dos Estudos 1 e 2}. Avaliar
capacidade de validação sobre rótulos internamente incoerentes mediria a incoerência,
não a capacidade.

\section{Estudo 1: validação}

\subsection{Uma tarefa diferente}

Aqui o agente não classifica. Ele recebe a decisão de um classificador numérico e julga
se ela é válida, dadas as estatísticas.

A tarefa é intrinsecamente mais fácil --- verificar é mais fácil que gerar --- e
operacionalmente mais valiosa: um sistema que sinaliza quando o classificador
provavelmente errou permite ao operador calibrar sua confiança caso a caso.

O conjunto tem 187 casos, sendo 149 decisões corretas e 38 incorretas, distribuídos
pelas sete classes não sintomáticas (1, 2, 3, 6, 7, 8 e 9).

\subsection{Resultados}

\begin{table}[htbp]
\centering
\caption{Desempenho do agente na tarefa de validação.}
\label{tab:ag-validacao}
\small
\begin{tabular}{lcc}
\toprule
\textbf{Métrica} & \textbf{Valor} & \textbf{IC $95\,\%$} \\
\midrule
Acurácia top-2 global                & $71{,}7\,\%$ ($134/187$) & $[64{,}8;\ 77{,}6]$ \\
\quad confirmação de decisões corretas & $71{,}8\,\%$ ($107/149$) & --- \\
\quad rejeição de decisões incorretas  & $71{,}1\,\%$ ($27/38$)   & --- \\
\midrule
Corretas julgadas VÁLIDAS            & $67{,}1\,\%$ ($100/149$) & --- \\
Incorretas julgadas INVÁLIDAS        & $73{,}7\,\%$ ($28/38$)   & --- \\
\midrule
\textbf{Precisão}                    & $\bm{0{,}909}$ ($100/110$) & $[0{,}841;\ 0{,}950]$ \\
Revocação                            & $0{,}671$ & --- \\
Medida $F_1$                         & $0{,}772$ & --- \\
\bottomrule
\end{tabular}
\end{table}

\begin{resultadochave}
Quando o agente declara uma classificação VÁLIDA, ele está certo em $\bm{90{,}9\,\%}$ dos
casos. E ele rejeita $73{,}7\,\%$ das classificações efetivamente erradas.
\end{resultadochave}

\subsection{Por que a precisão é o número que importa}

A assimetria entre precisão ($0{,}909$) e revocação ($0{,}671$) define o uso operacional
correto.

\begin{description}[leftmargin=0pt,style=nextline,itemsep=0.4em]

\item[Um selo de VÁLIDO é confiável.] Nove em cada dez são corretos. O operador pode
tratá-lo como forte evidência a favor da classificação.

\item[A ausência do selo não é evidência contra.] Cerca de um terço das classificações
corretas não recebe confirmação. Um caso não confirmado exige exame; não deve ser
descartado.

\end{description}

\begin{destaque}[O uso correto]
Este perfil sustenta uma política concreta de triagem. Segmentos com selo de VÁLIDO
podem ser processados com menor escrutínio; segmentos sem selo entram na fila de revisão
humana.

Como o agente confirma cerca de $59\,\%$ do total ($110$ de $187$), essa política reduz
em mais da metade o volume que exige atenção detalhada, com risco controlado de
aproximadamente $9\,\%$ de confirmações incorretas. É uma proposta de valor mensurável,
e é diferente --- e mais defensável --- de ``a IA classifica os eventos''.
\end{destaque}

\subsection{Por classe}

\begin{table}[htbp]
\centering
\caption{Desempenho do agente na validação, por classe (\%). ``Val.\ OK'' é a fração de
decisões corretas julgadas válidas; ``Val.\ Rej.'' é a fração de decisões incorretas
julgadas inválidas.}
\label{tab:ag-validacao-classe}
\small
\begin{tabular}{lccccc}
\toprule
\textbf{Cl.} & \textbf{Conf.\ top-1} & \textbf{Conf.\ top-2} &
\textbf{Erradas rejeitadas} & \textbf{Val.\ OK} & \textbf{Val.\ Rej.} \\
\midrule
1 & $40$ & $53$ & $60$ & $53$ & $60$ \\
2 & $95$ & $95$ & $80$ & $95$ & $80$ \\
3 & $62$ & $62$ & $40$ & $59$ & $60$ \\
6 & $73$ & $73$ & $80$ & $73$ & $80$ \\
7 & $86$ & $89$ & $\bm{100}$ & $81$ & $\bm{100}$ \\
8 & $47$ & $60$ & $40$ & $53$ & $60$ \\
9 & $29$ & $43$ & $88$ & $29$ & $75$ \\
\bottomrule
\end{tabular}
\end{table}

Duas linhas merecem comentário.

\paragraph{Classe 7.} Rejeita $100\,\%$ das classificações erradas e confirma $81\,\%$
das corretas. É o melhor validador da tabela --- notavelmente, para uma classe cuja
detecção de novidade atingia apenas $52{,}3\,\%$ no \cref{cap:mahalanobis}. As
competências não são as mesmas.

\paragraph{Classe 9.} Confirma apenas $29\,\%$ das corretas, mas rejeita $88\,\%$ das
erradas. O agente é um péssimo confirmador e um excelente cético para essa classe.

\begin{destaque}[Especialização por classe, medida e não postulada]
As \cref{tab:ag-porclasse,tab:ag-validacao-classe} juntas permitem atribuir a cada classe
o papel em que o agente é efetivamente bom:

classes 2 e 7 --- confirmar, justificar e validar;

classes 1, 3, 6 e 8 --- segunda opinião e nomeação de novidades;

classe 9 --- cético, acionado apenas para rejeitar;

classes 4 e 5 --- sinalizador sintomático, com reordenação das três candidatas para o
operador.

Essa atribuição não foi decidida a priori. Ela foi extraída dos dados, e o
\cref{cap:sintese} a incorporará à arquitetura de referência.
\end{destaque}

\section{Estudo 2: detecção de novidade}

\subsection{Protocolo}

Para cada classe, o agente recebe os perfis das \emph{demais} classes --- nunca o da
classe do segmento --- e deve reconhecer que nenhuma delas se aplica, além de propor um
nome e uma descrição. São 611 segmentos.

\subsection{Resultados}

\begin{table}[htbp]
\centering
\caption{Taxa de reconhecimento de novidade pelo agente, por classe.}
\label{tab:ag-novidade}
\small
\begin{tabular}{lrc}
\toprule
\textbf{Classe ocultada} & \textbf{Acertos} & \textbf{Taxa} \\
\midrule
1 & $51/72$   & $70{,}8\,\%$ \\
2 & $99/99$   & $\bm{100{,}0\,\%}$ \\
3 & $103/123$ & $83{,}7\,\%$ \\
6 & $67/70$   & $95{,}7\,\%$ \\
7 & $166/175$ & $94{,}9\,\%$ \\
8 & $56/61$   & $91{,}8\,\%$ \\
9 & $6/11$    & $54{,}5\,\%$ \\
\midrule
\textbf{Global} & $\bm{548/611}$ & $\bm{89{,}7\,\%}$ \\
\bottomrule
\end{tabular}
\end{table}

Intervalo de confiança de Wilson a $95\,\%$: $[87{,}0;\ 91{,}9]$.

\begin{resultadochave}
O agente reconhece corretamente que um segmento não pertence a nenhuma das classes
apresentadas em $\bm{89{,}7\,\%}$ dos casos --- contra $35{,}1\,\%$ de acerto quando
precisa nomear a classe correta.
\end{resultadochave}

\begin{destaque}[A assimetria central deste livro]
$35{,}1\,\%$ para dizer o que algo é. $89{,}7\,\%$ para dizer que não é nada conhecido.

A mesma assimetria apareceu no \cref{cap:mundoaberto}, quando o peso ótimo da OCSVM na
decisão de rejeição ($0{,}74$) inverteu-se em relação às decisões de reconhecimento. Ela
reaparece aqui, em um sistema de natureza completamente diferente.

Isso sugere que se trata de uma propriedade do \emph{problema}, não dos métodos:
reconhecer incompatibilidade é estruturalmente mais fácil que estabelecer identidade.
Basta que uma evidência contradiga cada classe conhecida; para afirmar identidade, é
preciso que todas as evidências convirjam.
\end{destaque}

\section{Os nomes gerados}

Ao rejeitar, o agente propõe um nome. A \cref{tab:ag-nomes} apresenta os nomes
consolidados por classe.

\begin{table}[htbp]
\centering
\caption{Nomes consolidados propostos pelo agente para cada classe apresentada como
desconhecida.}
\label{tab:ag-nomes}
\small
\begin{tabular}{lp{8.2cm}}
\toprule
\textbf{Cl.} & \textbf{Nome consolidado} \\
\midrule
1 & \emph{Wellhead Thermal-Pressure Surge} (sem convergência) \\
2 & \emph{Downhole Isolation Event with Wellhead Depressurization} \\
3 & \emph{Production Rate Surge with Cyclic Pressure Decline} \\
6 & \emph{Downhole Telemetry Loss with Dual Pressure Surge} \\
7 & \emph{Downhole Telemetry Loss with Wellhead Pressure Buildup} \\
8 & \emph{Flow Restriction with Thermal Cooling and Telemetry Loss} \\
9 & Sem convergência (seis nomes distintos) \\
\bottomrule
\end{tabular}
\end{table}

\subsection{O acerto notável}

Para a classe 2, o agente --- que nunca viu o rótulo ``fechamento espúrio de DHSV'' ---
produziu \emph{Downhole Isolation Event with Wellhead Depressurization}.

Traduzindo: ``evento de isolamento de fundo com despressurização na cabeça''. É uma
descrição fisicamente correta do fechamento espúrio da DHSV. O agente não recuperou o
nome; ele o \emph{reconstruiu} a partir do comportamento dos sensores.

\begin{destaque}[Por que isso importa mais que a taxa de acerto]
Um sistema que rejeita produz um balde de segmentos sem nome --- foi a limitação com que
o \cref{cap:mundoaberto} terminou. Um sistema que rejeita \emph{e propõe uma descrição
fisicamente plausível} produz uma hipótese que o especialista pode confirmar ou corrigir
em minutos, em vez de horas de análise.

É a diferença entre ``grupo 11, $n = 340$'' e ``possível evento de isolamento de fundo
com despressurização na cabeça''. A segunda saída é acionável.
\end{destaque}

\subsection{A não convergência}

A classe 1 produziu \textbf{51 nomes únicos} em suas 51 rejeições. A análise dos termos
mostra dispersão: \emph{thermal} aparece 16 vezes, \emph{pressure} 13, \emph{wellhead}
10, \emph{surge} 10. O nome consolidado da \cref{tab:ag-nomes} é uma síntese posterior,
não um consenso do agente.

A classe 9 produziu seis nomes distintos em seis rejeições.

\begin{destaque}[A não convergência é um sinal útil]
Quando o agente converge para um nome, há estrutura reconhecível. Quando cada segmento
gera um nome diferente, o agente está confabulando --- e isso pode ser \textbf{medido}
automaticamente, sem rótulos.

Propõe-se, portanto, a taxa de convergência de nomes como \emph{métrica de confiança não
supervisionada}: agrupar as descrições geradas e medir sua concentração. Alta
concentração indica descoberta genuína; alta dispersão indica que a saída não deve ser
apresentada ao operador como hipótese.
\end{destaque}

\subsection{O achado incômodo}

Quatro das sete classes --- 2, 6, 7 e 8 --- receberam nomes contendo
\emph{telemetry loss}, perda de telemetria.

Isso não é uma característica física dos eventos. É um artefato do 3W: em muitas
instâncias, o início da anomalia coincide com sensores de fundo reportando NaN ou zero.
O agente descreveu corretamente o que \emph{está nos dados}, e o que está nos dados
inclui uma correlação espúria entre anomalia e falha de telemetria.

\begin{destaque}[Quando o modelo revela o conjunto de dados]
O \cref{cap:dados} já havia alertado, na \cref{sec:dados-telemetria}, para essa
correlação. O que este resultado acrescenta é que ela é \textbf{forte o bastante para
dominar a descrição gerada} de quatro classes distintas.

A implicação é séria: qualquer modelo treinado no 3W pode estar aprendendo, em parte, a
detectar falha de telemetria em vez de detectar a anomalia. Os resultados de todos os
capítulos anteriores deste livro devem ser lidos com essa ressalva.

E há uma lição metodológica geral: um sistema que descreve o que vê em linguagem natural
funciona como instrumento de auditoria do conjunto de dados. Nenhuma das matrizes de
confusão dos capítulos anteriores tornou esse artefato visível.
\end{destaque}

\section{Limitações}

\begin{enumerate}[leftmargin=*,itemsep=0.4em]
  \item \textbf{Ausência de estudos de ablação.} Não se avaliou a contribuição individual
        de cada métrica, nem se comparou com métodos de atribuição como SHAP ou
        gradientes integrados.
  \item \textbf{Ausência de ajuste fino.} Modelo de propósito geral, sem adaptação ao
        domínio.
  \item \textbf{Passagem única de inferência.} Sem agregação de múltiplas amostragens,
        que tipicamente melhora consistência.
  \item \textbf{Apenas dados públicos.} Somente o 3W, com todos os seus artefatos.
  \item \textbf{Confundimento de telemetria.} Discutido acima; afeta quatro das sete
        classes avaliadas.
  \item \textbf{Amostras pequenas por classe.} A classe 9 tem $n = 11$ no Estudo 2. Os
        intervalos de confiança correspondentes são largos.
  \item \textbf{Possível contaminação.} O 3W é público desde 2019.
\end{enumerate}

\begin{destaque}[Por que os resultados ainda sustentam a conclusão]
A objeção mais séria é a contaminação: o modelo poderia estar recuperando conhecimento
memorizado sobre o 3W, e não raciocinando.

Três observações argumentam contra essa explicação. Primeiro, se houvesse memorização, a
acurácia de classificação seria \emph{alta}, não $35{,}1\,\%$. Segundo, os erros seguem
padrões fisicamente coerentes --- a confusão entre as classes 8 e 9 corresponde a
mecanismos idênticos, o que memorização de rótulos não produziria. Terceiro, os nomes
gerados não reproduzem a nomenclatura do 3W; eles a reconstroem em vocabulário próprio.

A contaminação provavelmente existe e provavelmente não explica os resultados.
\end{destaque}

\resumocapitulo{%
Três estudos medem o agente. Na classificação, ele acerta $35{,}1\,\%$ na primeira
escolha e $63{,}9\,\%$ nas três primeiras --- três vezes o acaso, e muito abaixo dos
$87\,\%$ dos classificadores numéricos; a diferença entre os níveis mostra que ele
identifica o conjunto plausível de hipóteses e falha em ordená-lo, o que o caracteriza
como gerador de hipóteses. Na validação, obtém precisão de $0{,}909$ com revocação de
$0{,}671$: um selo de VÁLIDO é confiável, sua ausência não é evidência contra, e isso
sustenta uma política de triagem que reduz à metade o volume de revisão humana. Na
detecção de novidade, atinge $89{,}7\,\%$ --- a assimetria entre $35{,}1\,\%$ para
afirmar identidade e $89{,}7\,\%$ para reconhecer incompatibilidade reaparece aqui em um
sistema completamente distinto daquele do \cref{cap:mundoaberto}, sugerindo tratar-se de
propriedade do problema. Os nomes gerados convergem para descrições fisicamente corretas
em cinco das sete classes, com destaque para a classe 2 --- \emph{Downhole Isolation
Event with Wellhead Depressurization} ---, e a taxa de convergência é proposta como
métrica de confiança não supervisionada. Quatro classes receberam nomes contendo
``perda de telemetria'', revelando um artefato do 3W forte o bastante para dominar a
descrição e exigir ressalva sobre todos os resultados do livro.}

\begin{exercicios}
\item Calcule os intervalos de confiança de Wilson para todas as taxas da
      \cref{tab:ag-novidade}. Para quais classes o intervalo é largo o bastante para
      tornar a comparação entre elas inconclusiva?

\item O agente acerta $35{,}1\,\%$ em top-1 e $63{,}9\,\%$ em top-3. Modele o processo
      como ordenação de nove candidatas e estime a distribuição da posição da resposta
      correta. Que fração da massa está além da terceira posição?

\item A precisão de $0{,}909$ com revocação de $0{,}671$ sustenta uma política de
      triagem. Formalize essa política, estime a redução de carga de revisão em função do
      volume diário de segmentos e calcule o risco residual.

\item Implemente a métrica de convergência de nomes proposta no capítulo: agrupe
      descrições geradas por similaridade semântica e meça a concentração. Valide-a
      verificando se ela distingue as classes 2 e 9.

\item A confusão entre as classes 8 e 9 é atribuída à ausência de informação de
      localização nas métricas. Proponha duas métricas adicionais que a codifiquem e
      preveja o efeito sobre a acurácia top-1 da classe 9.

\item Projete um experimento que quantifique o efeito do confundimento de telemetria:
      por exemplo, remover instâncias em que a anomalia coincide com valores ausentes e
      repetir a avaliação. Quanta acurácia você espera perder, e por quê?

\item Argumente a favor e contra a afirmação: ``a contaminação do conjunto de dados de
      treinamento invalida estes resultados''. Proponha um conjunto de avaliação que
      resolveria definitivamente a questão e discuta a viabilidade de construí-lo.
\end{exercicios}

\begin{leituras}
\item \citet{Lopes2026LLM} --- o trabalho que originou este capítulo, com os três estudos
      completos.
\item \citet{Wilson1927} --- o intervalo de confiança usado em todas as taxas.
\item \citet{VazVargas2026_3w2} --- a versão 2.0 do 3W, com discussão sobre qualidade dos
      dados relevante ao artefato de telemetria.
\item \citet{Sharma2025xaiagents} e \citet{Belay2026agenticAD} --- trabalhos
      contemporâneos sobre agentes e explicabilidade em detecção de anomalias.
\end{leituras}
